\section{随机数生成器}
\label{sec:random}

\subsection{背景}

游戏中的随机数无处不在：地形生成决定生物群系分布和矿石位置，生物AI决定行为选择，掉落系统计算掉落物品，天气系统决定雷电位置。随机数生成器的性能和质量直接影响游戏体验。

\subsection{原版局限性}

Java版Minecraft使用\texttt{java.util.Random}，基于48位线性同余生成器（LCG）：
\begin{equation}
X_{n+1} = (a \cdot X_n + c) \mod m
\end{equation}
其中$a = 25214903917$，$c = 11$，$m = 2^{48}$。主要问题包括：周期有限（$2^{48}$）、低维分布不均匀、线程安全锁开销、虚函数调用开销。

\subsection{本项目实现}

\subsubsection{xoroshiro128++算法}

默认采用xoroshiro128++算法，状态空间128位，周期$2^{128}-1$。核心算法包含两个步骤：

\textbf{输出计算}：$result = rotl(s_0 + s_1, 17) + s_0$

\textbf{状态更新}：$s_1 \leftarrow s_1 \oplus s_0$，然后$s_0$和$s_1$通过旋转和异或更新。

其中$rotl(x, k) = (x \ll k) | (x \gg (64-k))$是循环左移，现代CPU有专用指令支持。

种子初始化使用SplitMix64算法从单个64位种子生成两个64位状态字，确保不同种子产生差异大的初始状态。

\subsubsection{无偏差范围生成}

\texttt{nextInt(bound)}需要特别注意：简单的取模方法会产生偏差，因为$2^{32}$不是任意bound的倍数。解决方案是拒绝采样：当随机值超过阈值$\lfloor 2^{32}/bound \rfloor \cdot bound$时重新采样，确保每个结果出现概率严格相等。

\subsubsection{高斯分布}

使用Marsaglia极坐标法生成正态分布随机数。每次调用产生两个高斯值，缓存第二个值供下次使用，避免重复计算。

\subsubsection{算法选择}

通过编译宏可以选择不同算法：xoroshiro128++（默认，16字节状态，2-3倍速度提升）、xoshiro256++（32字节状态，极高质量）、LCG（8字节状态，最快）、MT19937（2.5KB状态，标准兼容）。

\subsection{成果}

\begin{table}[h]
\centering
\caption{随机数生成器对比}
\label{tab:random_comparison}
\begin{tabularx}{\linewidth}{Xcccc}
\toprule
\textbf{算法} & \textbf{周期} & \textbf{状态} & \textbf{速度} & \textbf{质量} \\
\midrule
Java LCG & $2^{48}$ & 8B & 1.0$\times$ & 一般 \\
xoroshiro128++ & $2^{128}-1$ & 16B & 2-3$\times$ & 高 \\
xoshiro256++ & $2^{256}-1$ & 32B & 2$\times$ & 极高 \\
\bottomrule
\end{tabularx}
\end{table}

xoroshiro128++相比Java LCG：周期更长、速度更快、状态更小、质量更高。配合无偏差的范围生成和高斯分布缓存，满足游戏各系统的需求。
